% ------------------------------------------------------------------------------
% Chapter 1
% ------------------------------------------------------------------------------
\chapter{Design Overview} 
\label{cha:design_overview}

The \texttt{gcd} module computes the Greatest Common Divisor of two unsigned integers using the subtractive Euclidean algorithm. It processes one input pair at a time and communicates through two independent ready/valid handshake channels.

\subsection*{Mathematical Definition}
For unsigned operands $A$ and $B$:
\begin{align*}
    \gcd(A, 0) &= A \\
    \gcd(0, B) &= B \\
    \gcd(0, 0) &= 0 \quad \text{(covers both above)} \\
    \gcd(A, A) &= A \\
    \gcd(A, B) &= \gcd(A - B, B) \quad \text{when } A > B \\
    \gcd(A, B) &= \gcd(A, B - A) \quad \text{when } B > A
\end{align*}

\subsection*{Operation Sequence}
\begin{enumerate}
    \item \textbf{IDLE:} The module waits in the \texttt{IDLE} state, asserting \texttt{in\_ready = 1}.
    \item \textbf{Handshake \& RUN:} On the input handshake, operands are latched. The result is determined immediately for zero or equal operands; otherwise, iterative subtraction begins in the \texttt{RUN} state.
    \item \textbf{DONE:} When convergence is reached, the result is latched into the output register and the module enters the \texttt{DONE} state, asserting \texttt{out\_valid = 1}.
    \item \textbf{Return:} After the output handshake completes, the module returns to \texttt{IDLE}.
\end{enumerate}

\section{Parameters}
\begin{table}[H]
    \centering
    \begin{tabular}{llp{8cm}}
        \toprule
        \textbf{Parameter} & \textbf{Default} & \textbf{Description} \\
        \midrule
        \texttt{WIDTH} & 16 (must be $\geq 1$) & Width of the data operands \\
        \bottomrule
    \end{tabular}
    \caption{Module Parameters}
\end{table}

\section{Interface}

\subsection*{Clock and Reset}
\begin{table}[H]
    \centering
    \begin{tabularx}{\textwidth}{l l c X}
        \toprule
        \textbf{Signal} & \textbf{Direction} & \textbf{Width} & \textbf{Description} \\
        \midrule
        \texttt{clk} & Input & 1 & Clock signal \\
        \texttt{rst\_n} & Input & 1 & Active-low sync reset \\
        \bottomrule
    \end{tabularx}
\end{table}

\subsection*{Input Channel}
\begin{table}[H]
    \centering
    \begin{tabularx}{\textwidth}{l l c X}
        \toprule
        \textbf{Signal} & \textbf{Direction} & \textbf{Width} & \textbf{Description} \\
        \midrule
        \texttt{in\_valid} & Input & 1 & Asserted by the producer when \texttt{a\_in} and \texttt{b\_in} hold valid data \\
        \texttt{in\_ready} & Output & 1 & Asserted by the module when it can accept a new transaction. High only in the \texttt{IDLE} state. \\
        \texttt{a\_in} & Input & \texttt{WIDTH} & First operand \\
        \texttt{b\_in} & Input & \texttt{WIDTH} & Second operand \\
        \bottomrule
    \end{tabularx}
\end{table}

\subsection*{Output Channel}
\begin{table}[H]
    \centering
    \begin{tabularx}{\textwidth}{l l c X}
        \toprule
        \textbf{Signal} & \textbf{Direction} & \textbf{Width} & \textbf{Description} \\
        \midrule
        \texttt{out\_valid} & Output & 1 & Asserted by the module when \texttt{gcd\_out} holds a valid result. High for the entire \texttt{DONE} state. \\
        \texttt{out\_ready} & Input & 1 & Asserted by the consumer when it is ready to accept the result. \\
        \texttt{gcd\_out} & Output & \texttt{WIDTH} & GCD result. Valid and stable for the entire period that \texttt{out\_valid} is asserted. \\
        \bottomrule
    \end{tabularx}
\end{table}

\section{Functional Description}
The behavior of the \texttt{gcd} module is driven by a Finite State Machine (FSM) comprising three states:

\subsection*{FSM States}
\begin{table}[H]
    \centering
    \begin{tabularx}{\textwidth}{l c c X}
        \toprule
        \textbf{State} & \textbf{\texttt{in\_ready}} & \textbf{\texttt{out\_valid}} & \textbf{Description} \\
        \midrule
        \texttt{IDLE} & 1 & 0 & Waiting for input. Ready to accept a new transaction. \\
        \texttt{RUN}  & 0 & 0 & Iterative subtraction in progress. No new input accepted. \\
        \texttt{DONE} & 0 & 1 & Result ready. Waiting for the consumer to accept it. \\
        \bottomrule
    \end{tabularx}
\end{table}

\subsection*{Algorithm Execution (\texttt{RUN} State)}
Each clock cycle while in the \texttt{RUN} state, the working registers are updated combinatorially for the next cycle:
\begin{itemize}
    \item \textbf{If} \texttt{a\_reg > b\_reg}: \texttt{a\_next = a\_reg - b\_reg} and \texttt{b\_next = b\_reg}.
    \item \textbf{Else if} \texttt{b\_reg > a\_reg}: \texttt{b\_next = b\_reg - a\_reg} and \texttt{a\_next = a\_reg}.
\end{itemize}

\subsection*{State Transitions}
\begin{description}
    \item[\texttt{IDLE} $\rightarrow$ \texttt{DONE}:] Triggered by an input handshake \textbf{AND} (\texttt{a\_in == 0 OR b\_in == 0 OR a\_in == b\_in}). The result is known immediately, making this a 1-cycle transition.
    \item[\texttt{IDLE} $\rightarrow$ \texttt{RUN}:] Triggered by an input handshake \textbf{AND} (\texttt{a\_in != 0 AND b\_in != 0 AND a\_in != b\_in}).
    \item[\texttt{RUN} $\rightarrow$ \texttt{DONE}:] Triggered when the next-cycle values of the working registers are equal (\texttt{a\_next == b\_next}). This evaluates against the combinatorial \textit{next-cycle} values, firing on the same clock edge as the final subtraction step.
    \item[\texttt{DONE} $\rightarrow$ \texttt{IDLE}:] Triggered by a successful output handshake (\texttt{out\_valid == 1 AND out\_ready == 1}).
\end{description}